% =====================================================================
%  chunks/refs.tex —— 参考文献（原书 p.589–591）
%  按本库房规：文献条目照排英文原文，不译；仅清理 PDF 文字层的重音与断字。
% =====================================================================
\chapter*{参考文献}
\addcontentsline{toc}{chapter}{参考文献}
\markboth{参考文献}{参考文献}

% [原文 p.589]
\begin{biblist}
\bibentry{1}{J.~Backelin, ``On the rates of growth of the homologies of Veronese subrings'',
Algebra, Algebraic Topology and their Interactions (Stockholm, 1983), Springer-Verlag
(Lecture Notes in Math., vol.~1183), Berlin 1985, pp.~79--100.}
\bibentry{2}{A.~A.~Beilinson, ``Coherent sheaves on $\PP^n$ and problems in linear algebra'',
Funktsional. Anal. i Prilozhen. \textbf{12}:3 (1978), 68--69; English transl.,
Functional Anal. Appl. \textbf{12} (1979), 214--216.}
\bibentry{3}{A.~Beilinson, J.~Bernstein, and P.~Deligne, ``Faisceaux pervers'', Analysis and
Topology on Singular Spaces, I (Luminy, 1981), Soc. Math. France (Ast\'erisque, vol.~100),
Paris 1982, pp.~5--171.}
\bibentry{4}{I.~N.~Bernshtein, I.~M.~Gel'fand, and S.~I.~Gel'fand, ``Algebraic vector bundles
on $\PP^n$ and problems of linear algebra'', Funktsional. Anal. i Prilozhen. \textbf{12}:3
(1978), 66--67; English transl., Functional Anal. Appl. \textbf{12} (1979), 212--214.}
\bibentry{5}{A.~I.~Bondal, ``Representations of associative algebras and coherent sheaves'',
Izv. Akad. Nauk SSSR Ser. Mat. \textbf{53} (1989), 25--44; English transl.,
Math. USSR-Izv. \textbf{34} (1990), 23--42.}
\bibentry{6}{A.~I.~Bondal and M.~Kapranov, ``Representable functors, Serre functors, and
reconstructions'', Izv. Akad. Nauk SSSR Ser. Mat. \textbf{53} (1989), 1183--1205;
English transl., Math. USSR-Izv. \textbf{35} (1990), 519--541.}
\bibentry{7}{A.~Bondal and D.~Orlov, ``Semi-orthogonal decomposition for algebraic varieties'',
Preprint alg-geom/9506006, Max-Planck-Institut f\"ur Mathematik, Bonn 1995.}
\bibentry{8}{A.~Bondal and D.~Orlov, ``Reconstruction of a variety from the derived category
and groups of autoequivalences'', Compositio Math. \textbf{125} (2001), 327--344.}
\bibentry{9}{T.~Bridgeland, ``Equivalences of triangulated categories and Fourier--Mukai
transforms'', Bull. London Math. Soc. \textbf{31}:1 (1999), 25--34.}
\bibentry{10}{T.~Bridgeland, ``Flops and derived categories'', Invent. Math. \textbf{147}
(2002), 613--632.}
\bibentry{11}{H.~Cartan and S.~Eilenberg, \emph{Homological algebra}, Princeton Univ. Press,
Princeton, NJ 1956; Russian transl., Izdat. Inostr. Lit., Moscow 1960.}
\bibentry{12}{P.~Deligne, ``Cohomologie \`a supports propres'', Th\'eorie des Topos et
Cohomologie \'Etale des Sch\'emas (SGA~4), Springer-Verlag (Lecture Notes in Math.,
vol.~305), Berlin 1973, pp.~252--480.}
\bibentry{13}{P.~Gabriel and M.~Zisman, \emph{Calculus of fractions and homotopy theory},
Springer-Verlag, New York 1967; Russian transl., Mir, Moscow 1971.}
\bibentry{14}{S.~I.~Gel'fand and Yu.~I.~Manin, \emph{Methods of homological algebra}, vol.~1,
Introduction to the theory of cohomology, and derived categories, Nauka, Moscow 1988;
English transl., Springer-Verlag, Berlin 1996.}
\bibentry{15}{A.~Grothendieck, ``Sur quelques points d'alg\`ebre homologique'', T\^ohoku
Math. J. (2) \textbf{9} (1957), 119--221.}
\bibentry{16}{A.~Grothendieck, ``\'El\'ements de g\'eom\'etrie alg\'ebrique'', Publ. Math.
Inst. Hautes \'Etudes Sci. \textbf{4}, 8, 11, 17, 20, 24, 28, 32 (1961--1967).}
\bibentry{17}{R.~Hartshorne, \emph{Residues and duality}, Springer-Verlag (Lecture Notes in
Math., vol.~20), Berlin 1966.}
\bibentry{18}{S.~P.~Inamdar and V.~B.~Mehta, ``Frobenius splitting of Schubert varieties and
linear syzygies'', Amer. J. Math. \textbf{116} (1994), 1569--1586.}

% [原文 p.590]
\bibentry{19}{V.~A.~Iskovskikh, \emph{Lectures on three-dimensional algebraic varieties:
Fano varieties}, Moscow State Univ., Moscow 1988. (Russian)}
\bibentry{20}{M.~M.~Kapranov, ``Derived category of coherent sheaves on Grassmann manifolds'',
Izv. Akad. Nauk SSSR Ser. Mat. \textbf{48} (1984), 192--202; English transl.,
Math. USSR-Izv. \textbf{24} (1985), 183--192.}
\bibentry{21}{M.~M.~Kapranov, ``Derived category of coherent bundles on a quadric'',
Funktsional. Anal. i Prilozhen. \textbf{20}:2 (1986), 67; English transl.,
Functional Anal. Appl. \textbf{20} (1986), 141--142.}
\bibentry{22}{M.~Kapranov, ``On the derived categories of coherent sheaves on some homogeneous
spaces'', Invent. Math. \textbf{92} (1988), 479--508.}
\bibentry{23}{M.~Kashiwara, ``The Riemann--Hilbert problem for holonomic systems'',
Publ. Res. Inst. Math. Sci. \textbf{20} (1984), 319--365.}
\bibentry{24}{M.~Kashiwara and P.~Schapira, \emph{Sheaves on manifolds}, Springer-Verlag,
Berlin 1990; Russian transl., Mir, Moscow 1994.}
\bibentry{25}{B.~Keller, ``Derived categories and their uses'', Handbook of Algebra, vol.~I
(M.~Hazewinkel, ed.), North-Holland, Amsterdam 1996, pp.~671--701.}
\bibentry{26}{H.~W.~Lenstra, Jr, F.~Oort, and Yu.~G.~Zarhin, ``Abelian subvarieties'',
J. Algebra \textbf{180} (1996), 513--516.}
\bibentry{27}{E.~Looijenga and C.~Peters, ``Torelli theorems for K\"ahler K3 surfaces'',
Compositio Math. \textbf{42}:2 (1980/81), 145--186.}
\bibentry{28}{M.~Miyanishi, ``Some remarks on algebraic homogeneous vector bundles'',
Number Theory, Algebraic Geometry and Commutative Algebra, Kinokuniya, Tokyo 1973,
pp.~71--93.}
\bibentry{29}{S.~Mukai, ``Duality between $D(X)$ and $D(\hat X)$ with its application to
Picard sheaves'', Nagoya Math. J. \textbf{81} (1981), 153--175.}
\bibentry{30}{S.~Mukai, ``Semi-homogeneous vector bundles on an Abelian variety'',
J. Math. Kyoto Univ. \textbf{18} (1978), 239--272.}
\bibentry{31}{S.~Mukai, ``On the moduli space of bundles on K3 surfaces. I'', Vector Bundles
on Algebraic Varieties (Bombay, 1984), Tata Inst. Fund. Res. (Tata Inst. Stud. Math.,
vol.~11), Bombay 1987, pp.~341--413.}
\bibentry{32}{D.~Mumford, \emph{Abelian varieties}, Oxford Univ. Press, London 1970.}
\bibentry{33}{V.~V.~Nikulin, ``Integer symmetric bilinear forms and some of their geometric
applications'', Izv. Akad. Nauk SSSR Ser. Mat. \textbf{43} (1979), 111--177; English
transl., Math USSR-Izv. \textbf{14} (1980), 103--167.}
\bibentry{34}{D.~O.~Orlov, ``Projective bundles, monoidal transformations, and derived
categories of coherent sheaves'', Izv. Ross. Akad. Nauk Ser. Mat. \textbf{56} (1992),
852--862; English transl., Russian Acad. Sci. Izv. Math. \textbf{41} (1993), 133--141.}
\bibentry{35}{D.~O.~Orlov, ``Equivalences of derived categories and K3 surfaces'',
J. Math. Sci. (New York) \textbf{84} (Algebraic geometry, 7) (1997), 1361--1381.}
\bibentry{36}{D.~O.~Orlov, ``Quasi-coherent sheaves in commutative and non-commutative
geometry'', Izv. Ross. Akad. Nauk Ser. Mat. \textbf{67}:3 (2003), 119--138; English
transl., Izv. Math. \textbf{67} (2003), 535--554.}
\bibentry{37}{D.~O.~Orlov, ``Derived categories of coherent sheaves on Abelian varieties and
equivalences between them'', Izv. Ross. Akad. Nauk Ser. Mat. \textbf{66}:3 (2002),
131--158; English transl., Izv. Math. \textbf{66} (2002), 569--594.}
\bibentry{38}{A.~Polishchuk, ``Symplectic biextensions and a generalization of the
Fourier--Mukai transform'', Math. Res. Lett. \textbf{3} (1996), 813--828.}
\bibentry{39}{I.~I.~Pyatetskii-Shapiro and I.~R.~Shafarevich, ``Torelli's theorem for
algebraic surfaces of type K3'', Izv. Akad. Nauk SSSR Ser. Mat. \textbf{35} (1971),
530--572; English transl., Math. USSR-Izv. \textbf{5} (1971), 547--588.}
\bibentry{40}{R.~G.~Swan, ``Hochschild cohomology of quasi-projective schemes'',
J. Pure Appl. Algebra \textbf{110}:1 (1996), 57--80.}
\bibentry{41}{\emph{Th\'eorie des intersections et th\'eor\`eme de Riemann--Roch},
S\'eminaire de G\'eom\'etrie Alg\'ebrique du Bois-Marie 1966--1967 (SGA~6), dirig\'e par
P.~Berthelot, A.~Grothendieck, L.~Illusie, avec la collaboration de D.~Ferrand,
J.~P.~Jouanolou, O.~Jussila, S.~Kleiman, M.~Raynaud et J.~P.~Serre, Springer-Verlag
(Lecture Notes in Math., vol.~225), Berlin 1971.}

% [原文 p.591]
\bibentry{42}{\emph{Th\'eorie des topos et cohomologie \'etale des sch\'emas}, S\'eminaire de
G\'eom\'etrie Alg\'ebrique du Bois-Marie 1963--1964 (SGA~4), dirig\'e par M.~Artin,
A.~Grothendieck et J.~L.~Verdier, avec la collaboration de N.~Bourbaki, P.~Deligne et
B.~Saint-Donat, vols.~1--3, Springer-Verlag (Lecture Notes in Math., vols.~269, 270, 305),
Berlin 1972--1973.}
\bibentry{43}{R.~Thomason and T.~Trobaugh, ``Higher algebraic $K$-theory of schemes and of
derived categories'', Grothendieck Festschrift, vol.~III, Birkh\"auser, Boston 1990,
pp.~247--435.}
\bibentry{44}{J.-L.~Verdier, ``Cat\'egories d\'eriv\'ees'', S\'eminaire de G\'eom\'etrie
Alg\'ebrique du Bois--Marie (SGA~4$\frac{1}{2}$), Springer-Verlag (Lecture Notes in Math.,
vol.~569), Berlin 1977, pp.~262--311.}
\end{biblist}

\medskip
\noindent\footnotesize
Steklov Institute of Mathematics, Russian Academy of Sciences\\
E-mail address: \texttt{orlov@mi.ras.ru}\\
Received 05/FEB/03. Translated by IPS(DoM). Typeset by \textsc{ams}-\TeX.
\normalsize

% [Tutor Note: 原书 [6] 条目在文字层与页图上均印作 “Representablefunctors”，缺一词间空格；本译本按页图照录，未改。]
% [Tutor Note: 原书 [44] 的 SGA 4 1/2 的分式在文字层被抽坏（(SGA 4 1 2 )），此处按页图排作 SGA 4½。]
